\bookchapter{第八章：通信：短信、电话与社交}

\chapterintro{通信是手机最原始、最核心的使命。现在，它已从简单的通话和短信，演变成一个包含即时通讯、社交网络和电子邮件在内的复杂生态系统。}

\booksection{一、电话App的高级功能}

现代手机的电话应用，早已超越了一个简单的拨号盘，它集成了录音、骚扰拦截、智能识别等强大工具。

\booksubsection{通话录音（注意法律合规性）：}

\begin{itemize}
  \item \textbf{作用}：记录重要通话内容，用于工作纪要、采访、客户承诺或争议留存。
  \item \textbf{法律提示}：在中国境内，通话录音作为证据使用的前提是取得对方同意，或用于维护自身合法权益。
  \item \textbf{单方录音vs双方同意}：法律层面，"未经对方同意"取得的录音并非当然无效。最高人民法院相关司法解释中，只要录音取得方式未严重侵害他人合法权益、未违反法律禁止性规定、未违背公序良俗，仍可作为证据提交。但部分省份法院在民事审判中对"偷偷录音"的采信标准更严格；刑事案件则通常要求更严格的合法性审查。
  \item \textbf{各地差异}：北京、上海、广东等地法院对单方录音作为民事证据较为宽容；而涉及商业秘密、个人隐私的录音，无论在哪个省份都需要特别谨慎，可能因侵害隐私权而被排除。
  \item \textbf{如何操作}：安卓通常在通话界面有明确的"录音"按钮；iOS系统不提供原生通话录音功能，需借助第三方App或使用带录音功能的蓝牙耳机。
\end{itemize}

\begin{tcbwarn}{录音使用红线}
\begin{itemize}
  \item 不得将录音用于敲诈、传播或公开侮辱他人。
  \item 涉及第三方隐私（如偶然录到他人对话）的内容不得随意扩散。
  \item 重要场景建议在通话开始时口头告知"本次通话将被录音"，既合规也更具证据效力。
\end{itemize}

\end{tcbwarn}

\booksubsection{骚扰电话治理深度内容：}

\begin{itemize}
  \item \textbf{原理}：电话App通过网络查询庞大的号码库，在来电时显示"骚扰电话"、"诈骗电话"或公司名称。
  \item \textbf{如何开启}：安卓进入电话App > 设置 > 骚扰拦截；iOS进入设置 > 电话 > 来电阻止与身份识别。\textbf{各品牌详细开关路径与推荐开启项，见第十三章「实用系统工具」中的品牌对照表。}
  \item \textbf{主流标记平台}：国内号码标记主要由三家数据服务商提供——腾讯手机管家、360手机卫士、电话邦（Tapir）。运营商（移动/联通/电信）的"机主姓名"显示、"骚扰拦截"功能也常调用这些数据。手机厂商（华为、小米、OPPO、vivo）的标记库多与上述服务商合作。
  \item \textbf{被误标记怎么办}：如果你正常工作号码被错误标记为"外卖""快递""骚扰电话"，可通过以下渠道申诉：
  \begin{itemize}
    \item 电话邦：访问号码申诉平台（haoma.tk.cn）提交身份证明。
    \item 360：通过360手机卫士App内"号码申诉"功能提交。
    \item 腾讯：通过"腾讯手机管家"小程序"号码申诉"页面提交。
    \item 三大运营商：中国移动用户可发送短信"KTFSR"到10086开通防骚扰功能；电信用户关注"天翼防骚扰"公众号；联通用户在"联通手机管家"中设置。
  \end{itemize}
  \item \textbf{防骚扰进阶}：开启运营商的"高频呼叫拦截"（如移动的高频呼叫限制），可自动拦截AI语音外呼。
  \item \textbf{手机自带拦截常见能力}：
  \begin{itemize}
    \item \textbf{来电识别 / 智能拦截}：按号码库自动标记并拦截骚扰、诈骗、广告推销。
    \item \textbf{黑名单 / 白名单}：手动拉黑已知骚扰号；把通讯录、快递外卖号加入白名单，避免误杀。
    \item \textbf{短信关键词过滤}：包含"中奖""贷款""点击链接"等关键词的短信自动归入拦截箱。
    \item \textbf{iOS 静音未知来电}：不在通讯录中的来电只静音进语音信箱，不打断你；适合骚扰电话多的用户，但要留意漏接陌生重要电话。
    \item \textbf{勿扰模式联动}：夜间或开会时只允许通讯录来电，配合定时开关，睡眠质量立刻改善。
  \end{itemize}
  \item \textbf{实用软件怎么配合}：优先用系统自带拦截，再叠加国家反诈中心与运营商防骚扰；清理类可选清浊等轻量工具，但同一类安全/清理软件装一套即可，详见第十三章。
\end{itemize}

\booksubsection{VoLTE/VoNR是什么及如何开启：}

\begin{itemize}
  \item \textbf{VoLTE（Voice over LTE）}：基于4G网络的语音通话技术，通话时不再回落到2G/3G，能边打电话边上网，音质比传统通话更清晰，接通更快。
  \item \textbf{VoNR（Voice over NR）}：基于5G的语音通话，是VoLTE的下一代演进。延迟更低、音质更好，并能支持高清视频通话。需5G手机+5G套餐+当地5G覆盖三者具备才能使用。
  \item \textbf{开启方法}：iOS在"设置 > 蜂窝网络 > 蜂窝数据选项 > 语音与数据"中开启"VoLTE"；安卓在"设置 > 移动网络 > 移动数据 > 高级"中找到"VoLTE高清通话"开关。运营商通常默认开启，若发现通话时网络断开、无法上网，多半是VoLTE未启用。
\end{itemize}

\begin{tcbinfo}{小知识}
如果你的手机在通话时上网速度突然变慢或断网，多半是VoLTE/VoNR未开启，回落到了2G/3G网络。开启后即可实现“边打电话边刷视频”。

\end{tcbinfo}

\booksection{二、短信与iMessage/RCS}

传统短信并未被淘汰，它仍是接收验证码、银行通知和官方信息的可靠渠道。

\begin{itemize}
  \item \textbf{iMessage}：苹果设备间通过互联网发送的消息，显示为蓝色气泡。支持高质量图片视频、已读回执等，完全免费且端到端加密。
  \item \textbf{RCS}：由全球运营商和谷歌推动的下一代通信协议。功能类似iMessage或微信，但基于电话号码。
  \item \textbf{iMessage与RCS的互操作性}：iOS 18起，苹果开始支持RCS标准，这意味着iPhone与安卓手机之间可通过RCS发送高清图片、已读回执和打字状态提示，告别"绿色气泡"的体验断层。但前提是双方运营商均开通RCS，且手机系统版本足够新。
\end{itemize}

\booksubsection{RCS在国内的部署情况：}

\begin{itemize}
  \item \textbf{中国移动}：较早部署RCS，部分省份已支持"5G消息"业务，可发送富媒体消息、互动卡片。需手机终端支持且在营业厅开通。
  \item \textbf{中国电信}：通过"电信5G消息"提供服务，部分套餐默认支持，老用户需发送指定短信指令开通。
  \item \textbf{中国联通}：在主要城市部署5G消息，但终端覆盖参差不齐，许多安卓机型仍需手动在短信App中开启"增强信息"或"5G消息"开关。
  \item \textbf{实际体验}：国内RCS推广速度慢于预期，主要受微信生态强势、终端适配不统一等因素影响。日常通信仍以微信为主，RCS更多承担"官方推送"角色（如银行通知、政务消息）。
\end{itemize}

\begin{tcbwarn}{验证码安全红线}
\begin{itemize}
  \item \textbf{不要转发验证码}：任何让你"读出验证码"或"转发短信"的人都可能是诈骗分子，包括自称"客服""快递员""公检法"的人。验证码 = 你的钱袋子。
  \item \textbf{防钓鱼短信}：警惕"ETC过期""医保卡冻结""快递丢失理赔"等短信中的链接，正规机构不会通过短信链接索要银行卡号、密码。
  \item \textbf{号码识别}：银行官方短信通常来自95开头短号或1069+企业签名，凡是从普通11位手机号发来的"银行通知"基本都是诈骗。
  \item \textbf{双因子备份}：绑定重要账号的手机号不要随意更换；如必须换号，务必先解绑所有关键账号，否则新机主可能通过短信验证码重置你的密码。
\end{itemize}

\end{tcbwarn}

\booksection{三、微信/QQ：国民级社交应用的管理}

微信和QQ已超越通信工具，成为集社交、支付、信息获取于一体的“超级App”。

\booksubsection{聊天记录备份与迁移：}

\begin{enumerate}
  \item 在电脑上登录PC版或Mac版微信。
  \item 点击左下角菜单 > [迁移与备份] > [备份与恢复]。
  \item 选择"备份聊天记录至电脑"，在手机上选择需要备份的聊天。
  \item 迁移到新手机：在原手机微信[我] > [设置] > [通用] > [聊天记录迁移与备份] > [迁移] > [迁移到手机/平板]，扫码即可完成。
\end{enumerate}

\booksubsection{存储空间清理深度方法：}

\begin{itemize}
  \item \textbf{常规入口}：微信[我] > [设置] > [通用] > [存储空间]，可清理缓存和按时间、文件类型精细化清理聊天记录。
  \item \textbf{聊天内清理}：进入某个聊天 > 右上角"…" > [聊天记录] > [管理存储空间]，可单独清理某聊天中的图片、视频、文件，保留文字。
  \item \textbf{原图/原视频清理}：群聊中接收的高清原图视频最占空间，可在[管理存储空间]中勾选"原图"批量删除。
  \item \textbf{关闭自动下载}：[设置] > [通用] > [照片、视频、文件和通话]，关闭"自动下载"，避免在WiFi下默默缓存大量群文件。
  \item \textbf{小程序缓存}：在[发现] > [小程序]中长按不常用的小程序删除，可释放其本地缓存。
  \item \textbf{聊天记录压缩备份}：电脑版微信备份时勾选"压缩"，可将几GB的聊天记录压缩到原体积的30\%-50\%。
\end{itemize}

\booksubsection{群管理技巧：}

\begin{itemize}
  \item \textbf{群公告}：群主/管理员在群聊右上角"…" > [群公告]中发布，全员强提醒，适合重要通知。新成员入群时自动弹出。
  \item \textbf{群待办}：群主可将置顶消息设为"群待办"，所有成员头顶会出现红色提示，完成点击后自动消失。适合活动报名、任务收集。
  \item \textbf{群接龙}：在输入框输入"\#接龙"可发起接龙，适合点餐、统计报名。接龙结果自动汇总到群聊顶部，避免刷屏。
  \item \textbf{群直播}：群聊右上角"+" > [群直播]，可发起仅供本群观看的直播，适合企业内训、家族聚会、产品内测。
  \item \textbf{折叠该群聊}：在群列表向左滑动群名 > [折叠]，可将不重要的群收纳到"折叠的群聊"中，仅在新@你时提醒。
  \item \textbf{仅群主/管理员@所有人}：群成员500人以上时，普通成员无法@所有人，避免被打扰。
\end{itemize}

\booksubsection{朋友圈权限精细化管理：}

\begin{itemize}
  \item \textbf{分组可见性}：发布时点击"谁可以看"，选择"部分可见"或"不给谁看"，基于预先设置的标签分组（如"同事""家人""客户"）精准推送。
  \item \textbf{允许朋友查看范围}：[我] > [设置] > [朋友权限] > [朋友圈] > [允许朋友查看朋友圈的范围]，可选"最近三天/一个月/半年/全部"，建议设为"半年"以平衡隐私与社交。
  \item \textbf{屏蔽某条朋友圈}：长按朋友圈某条动态 > [不看他/她的朋友圈]或[不再关注]，对方不会知道。
  \item \textbf{朋友圈广告管理}：[我] > [设置] > [关于微信] > [个人信息与权限] > [个性化广告管理]，可关闭基于行为的个性化推荐。
  \item \textbf{陌生人十条}：[朋友权限] > [朋友圈] > [允许陌生人查看十条朋友圈]，关闭可避免陌生人翻看你的过往动态。
\end{itemize}

\booksection{四、视频通话对比}

不同视频通话工具各有侧重，根据场景选择合适的工具能让沟通更顺畅。

\begin{itemize}
  \item \textbf{微信视频}：国民级覆盖，老少咸宜。适合亲友日常联系、快速沟通。最多支持20人视频（需小程序"群组视频"），缺点是网络差时画质下降明显。
  \item \textbf{FaceTime}：苹果生态专属，画质优秀、端到端加密。iOS 15起支持"链接"邀请安卓/Web用户加入，但主持人必须是苹果设备。适合家庭跨设备沟通。
  \item \textbf{Zoom}：国际会议首选，支持最多100人免费40分钟会议。屏幕共享、虚拟背景、分组讨论室等功能完善，适合跨国/远程协作。
  \item \textbf{腾讯会议}：国内办公首选，与微信深度整合，扫码即入。支持300人会议（付费版），云端录制、实时字幕、会议纪要功能强大，适合企业内训和大型会议。
  \item \textbf{使用场景建议}：1对1日常联系选微信；苹果用户之间选FaceTime；正式商务会议选Zoom或腾讯会议（国内选腾讯，跨境选Zoom）。
\end{itemize}

\booksubsection{网络带宽要求与流量消耗参考：}

\begin{itemize}
  \item \textbf{微信视频通话}：约需500Kbps-1Mbps带宽，标准画质每分钟约消耗3-5MB流量，高清约8-12MB。
  \item \textbf{FaceTime}：标清约1Mbps，720p约2Mbps，1080p需4Mbps以上。每分钟流量约5-15MB。
  \item \textbf{Zoom/腾讯会议}：720p约1.5-3Mbps，1080p需3-5Mbps。每小时流量约300-800MB，多人会议主持人流量更大。
  \item \textbf{建议}：长时间视频通话尽量使用WiFi；移动网络下可主动降低画质（在App设置中关闭"高清视频"）以节省流量。
\end{itemize}

\booksubsection{视频通话画质优化技巧：}

\begin{itemize}
  \item \textbf{光线}：背对窗户、面朝光源（自然光最佳），避免逆光导致人脸变黑。
  \item \textbf{角度}：手机摄像头略高于眼睛，避免双下巴；横屏通话时双手持握更稳定。
  \item \textbf{音频}：使用耳机可避免回声和啸叫；安静环境下不要开"扬声器"模式。
  \item \textbf{网络}：优先WiFi；4G/5G视频时尽量靠近窗户避免信号衰减；多人会议可关闭其他视频画面仅保留发言人。
  \item \textbf{虚化背景}：腾讯会议、Zoom均支持AI虚化，可遮蔽杂乱背景。注意穿戴纯色衣物，避免与背景色相近导致"绿幕穿帮"。
\end{itemize}

\begin{tcbinfo}{小贴士}
网络卡顿时，先关闭视频保留语音，比直接挂断更礼貌。重要会议前可用“腾讯会议”自带的“网络检测”功能提前测试带宽。

\end{tcbinfo}

\booksection{五、电子邮件：在手机上配置与使用}

\begin{itemize}
  \item \textbf{添加邮箱账户}：iOS在设置 > 邮件 > 账户 > 添加账户；安卓通常有"电子邮件"App。
  \item \textbf{使用技巧}：为不同邮箱账户设置不同的通知方式；善用滑动手势快速归档、删除或标记邮件。
\end{itemize}

\booksubsection{IMAP vs POP3 vs Exchange协议区别：}

\begin{itemize}
  \item \textbf{POP3}：从服务器下载邮件后默认删除原件，仅本地保留。适合只用一台设备、需要离线阅读的场景。缺点是换设备后看不到历史邮件，已不推荐。
  \item \textbf{IMAP}：邮件始终保存在服务器，本地仅作缓存。多设备（手机/电脑/平板）同步状态一致——一端已读，全部已读。是当前主流推荐协议。
  \item \textbf{Exchange（含ActiveSync）}：微软协议，除邮件外还同步日历、联系人、任务。企业邮箱首选。Outlook、网易企业邮箱、腾讯企业邮箱均支持。
  \item \textbf{选择建议}：个人邮箱用IMAP；公司邮箱优先Exchange，其次IMAP；尽量不要用POP3，除非有特殊归档需求。
\end{itemize}

\booksubsection{企业邮箱配置方法：}

\begin{enumerate}
  \item iOS：[设置] > [邮件] > [账户] > [添加账户] > 选择"Microsoft Exchange"或"其他"。
  \item 安卓：在自带的"电子邮件"App或Outlook App中选择"添加账户"，输入完整邮箱地址和密码。
  \item 服务器填写（以Exchange为例）：服务器地址通常为"mail.公司域名.com"或"outlook.office365.com"，域名填公司简称。
  \item 若提示"无法验证服务器"，可手动输入：传入服务器（IMAP）端口993，加密SSL；传出服务器（SMTP）端口465或587，加密TLS。
  \item 登录后建议在[账户设置]中开启"推送"（Push）以实时接收邮件，并关闭"邮件预览"以保护商业机密。
\end{enumerate}

\begin{tcbinfo}{小提示}
企业邮箱通常需要管理员开通“客户端授权密码”或“IMAP/SMTP授权码”，不能直接用网页登录密码。登录失败时先到网页版邮箱设置中生成授权码再填入手机。

\end{tcbinfo}

\begin{tcbpractice}{实操：备份一次微信聊天记录到电脑}
\begin{enumerate}
  \item 在电脑上启动备份：打开电脑版微信 > 左下角菜单 > [迁移与备份] > [备份与恢复]。
  \item 选择备份到电脑：点击"备份聊天记录至电脑"。
  \item 在手机上选择备份内容，等待备份完成。
\end{enumerate}

\end{tcbpractice}
